\section[拓扑场论中的谱分解]{拓扑场论中的谱分解（4 月 15 日）}

\srcpage{62}

回忆我们开始讨论局部算子，把它看作在一点处（或附近）测量场。局部算子上的主要结构是算子乘积展开，它是一个乘法。给定两个靠近的点 $x,y$，在拓扑场论中计算局部算子的办法是：在这两点周围各挖去一个小球面，得到从两个球面到空集的协边。但也可以取一个同时包住 $x$ 与 $y$ 的更大球面，得到同一个局部算子，用同样的方式计算。于是我们得到协边 $S^{n-1}\to\varnothing$，而局部算子定义了 $Z(S^{n-1})$ 中的一个态。

若理论不是共形的，这要麻烦一点。要定义单个局部算子，得把点周围的球面缩小，取半径趋于零的极限。要做算子乘积展开，还必须把 $x$ 与 $y$ 相互靠近。

回到拓扑情形。算子乘积展开完全发生在以较大球面为边界的球 $B$ 内。把 $B$ 内 $x$、$y$ 周围的两个小球挖掉；于是 $B$ 减去这两个球面给出一个协边
\[
S^{n-1}\amalg S^{n-1}\longrightarrow S^{n-1},
\]
TFT 把它送到
\[
*:Z(S^{n-1})\otimes Z(S^{n-1})\longrightarrow Z(S^{n-1}).
\]
这就是算子乘积展开映射，希望它看起来像一个乘法。更确切地说，对更大球面内任意一对球面都得到这样的映射。一般结构是一族映射，参数空间是更大空间内两个小 $(n-1)$-球面的构形空间（configuration space）。由于理论是拓扑的，这族映射必须是局部常值的。

\begin{example}
当 $n=1$ 时，我们看的是更大线段内的一对区间。这个构形空间不连通：$\pi_0$ 知道 $x$ 是否在 $y$ 左边。因此对一维 TFT，算子乘积展开一般并不交换：
\[
O_y*O'_x\neq O_x*O'_y .
\]
不过它结合：更大区间内三个区间、并指定顺序的构形空间是连通的。
\end{example}

当 $n>1$ 时，更大圆盘内一对圆盘的构形空间连通，因此 $Z(S^{n-1})$ 是结合交换代数。有足够的空间把算子挪来挪去——这在非拓扑 QFT 中很微妙，但在拓扑情形特别好。而无论哪种情形，这种交换性在量子力学（$n=1$，此时 OPE 只是结合）中完全不存在。

\begin{remark}
下次我们要在导出情形下工作，那时这个故事会变得更复杂。我们会在 $Z(S^{n-1})$ 上得到的交换性概念（导出版）中，看到关于构形空间的更多信息。
\end{remark}

给定交换代数，我们想对它作谱分解，所以令 $\cM^{\mathrm{aff}}:=\Spec(Z(S^{n-1}))$。它称为理论的 vacua 模空间（moduli of vacua）（TODO：我想是这样）。给定局部算子 $\cO$ 与一个「真空」$|0\rangle$（即 $\cM^{\mathrm{aff}}$ 上的一个点），可以定义 $\cO$ 的真空期望值（vacuum expectation value, VEV），记作 $\langle 0|\cO|0\rangle$：算子是 $\cM^{\mathrm{aff}}$ 上的函数。在 TFT 中，我们把这个作为真空的定义。

口号是：拓扑场论（或者说它的局部算子）在 $\cM^{\mathrm{aff}}$ 上谱分解。TFT 中能算的一切都是 $Z(S^{n-1})$-线性的，于是分解到它的谱上。

\srcpage{63}

而且是一切。例如设 $M$ 是 $(n-1)$-流形，我们要在 $Z(M)$ 上造一个 $Z(S^{n-1})$-作用。具体地，考虑柱体 $M\times[0,1]$，在 $t=1/2$ 处、$M$ 中某点的周围挖去一个小球。修改后的流形是协边 $M\amalg S^{n-1}\to M$，于是定义了映射 $Z(S^{n-1})\otimes Z(M)\to Z(M)$，这就是作用映射。它依赖于我们选的点，但同样是局部常值的，这对我们够用（同样，导出情形下更微妙）。另一个协边说明它与 OPE 相容，于是我们得到模结构。过渡到谱，$Z(M)$ 定义了 $\cM^{\mathrm{aff}}$ 上的拟凝聚层（quasicoherent sheaf）。

另一种看法在导出情形下特别有用，那里「局部常值」更重要：$Z(S^{n-1})$ 在 $M$ 的所有点上的张量积作用在 $Z(M)$ 上。不过这高度冗余，它通过一个商分解，记作
\[
\int_M Z(S^{n-1}),
\]
称为 $Z(S^{n-1})$ 在 $M$ 上的因子化同调（factorization homology）。我们稍后会回到它。

我们先在二维有限规范理论中看这种层化（sheafification）。心里的奖赏在四维：这个构造在某种意义上是几何 Langlands 对应的来源。

\begin{example}[二维 TFT]
设 $\cC:=Z(\mathrm{pt}_+)$，其中 $Z$ 是扩张 TFT，于是 $\cC$ 是 $Z$ 的边界条件范畴。设 $\End(\id_\cC)$ 是 Bernstein 中心；例如当 $\cC=\Rep G$ 时，它是类函数代数 $\C[G]^G$，其交换乘法不是逐点的，而是卷积。

那么对这个 TFT $Z_G$，$\cM^{\mathrm{aff}}$ 是 $\Spec Z_G(S^1)=\Spec\C[G/G]$，我们把它定义为 $G$ 的对偶 $\hat G$。当 $G$ 交换时，这实现 Pontrjagin 对偶。但对任意有限 $G$，$\hat G$ 与 $G$ 的不可约表示的同构类集合一一对应。
\end{example}

这恢复了你可能见过的想法：研究有限群的表示时，人们想研究它们如何分解为不可约表示，即在 $\hat G$ 上层化。事实上整个范畴都分裂：
\[
\Rep G\simeq\bigoplus_{\text{irreps }V_\lambda}\Vect\otimes V_\lambda .
\]
也就是说，$\Rep G\simeq\Vect_{\hat G}$。

现在讨论 $Z(S^1)$ 在态空间上的作用。没有更多数据的话这不太有趣；但通过在 $\Rep G$ 中选两个对象，我们可以为边界染成红色和蓝色的紧一维流形定义态空间。例如把 $[0,1]$ 的 $0$ 端染红、$1$ 端染蓝，相应的协边定义了 $Z(S^1)$ 在 $Z([0,1])=\Hom(V,W)$ 上的作用。\yizhu{原文此处是一幅协边图（公式编号 22.7），纯文本底本没有保留这幅图；该协边是正方形内挖去一个圆所得的带边曲面。}{A（原文 p.63）}也就是说，这个 TFT 中的态空间与态射都富化在 $Z(S^1)$-模中。

某种意义上，Bernstein 中心 $Z(\cC)$ 对这种结构是万有的。$\cC$ 构成 $\cM:=\Spec(Z(\cC))$ 上的范畴层。\footnote{我们为 TFT 使用记号 $Z$，不要与代数、范畴等的中心记号 $Z$ 混淆。}若 $R$ 是交换环，则 $\cC$ 上的 $R$-线性结构等价于环同态 $R\to Z(\cC)$ 的数据：$R$ 以相容的方式作用在 Hom 空间上。因此万有的这种 $R$ 就是 $Z(\cC)$。

如果你处在足够大的环境里（$\cC$ 有足够多的余极限），$R$-线性等价于说 $\cC$ 是 $\mathrm{Mod}_R$ 上的模范畴。我们现在正在做范畴的谱分解，尽管只在仿射空间上。

把范畴层化到 $\Spec R$ 上，有三种不同的定义方式。
\begin{enumerate}
\item 把 $\cC$ 在 $R$ 上层化，是给一个映射 $R\to Z(\cC)$ 的数据。
\item 把 $\cC$ 在 $R$ 上层化，意思是 $\cC$ 富化在 $R$-模中，等价地富化在 $\Spec R$ 上的拟凝聚层中。
\item 把 $\cC$ 在 $R$ 上层化，意思是令对称幺半范畴 $(\mathrm{Mod}_R,\otimes_R)$ 作用在 $\cC$ 上。
\end{enumerate}

\srcpage{64}

对最后一种视角，我们知道若 $M$ 是自由 $R$-模、$C\in\cC$，则 $M\otimes C\cong C^{\oplus\rank(M)}$。一般地，取 $M$ 的自由消解，于是得到 $\cC$ 中的一串对象。若 $\cC$ 有足够多的余极限，可以取这串的余极限，得到我们称为 $M\otimes C$ 的东西。这个作用是线性的范畴类比。

也可以局部化：若 $S\subset R$ 是乘性子集，我们希望能把 $\cC$ 局部化到 $\Spec R[S^{-1}]$ 上，就像可以把拟凝聚层局部化到 $\Spec R[S^{-1}]$ 上。定义层
\[
\cC(\Spec R[S^{-1}])=\cC\otimes_{\mathrm{Mod}_R}\mathrm{Mod}_{R[S^{-1}]}.
\]
具体地，这取 $\cC$ 并把来自 $S\subset R$ 的所有态射取逆。我们暂时不担心范畴张量积的精确定义。但看清一点也许有帮助：$\cC$ 中的 Hom 集已经是 $R$ 上的层，而局部化 $\cC$ 归结为对 Hom 集对应的层或模作局部化。

现在设 $A$ 是结合代数，例如 $\C[G]$。则 $\mathrm{Mod}_A$ 层化到 $\Spec Z(A)$ 上，特别地 $\Rep G$ 层化到 $\hat G$ 上。因此事实上 $\Rep G\simeq\QCoh(\hat G)$——但这太好了，不可能是真的！这是一个非典范的等价。

事实上你早该起疑：$\Rep G$ 在 $\hat G$ 上层化是同义反复。对 $\hat G$ 的子集 $S$，$\Rep G$ 在 $S$ 上的限制是特征标都落在 $S$ 中的那个直和项。让 $S$ 是一个点，就看出了问题所在：要得到 $\Rep G\simeq\QCoh(\hat G)$ 的等同，我们需要选取 $G$-表示的同构类。

换种说法，Wedderburn 定理给出
\[
\C[G]\simeq\bigoplus_{\text{irreps }V_\lambda}\End(V_\lambda),
\]
是除环上矩阵代数的直和，都是环上的代数。我们已经把 $\C[G]$ 分解为 $\hat G$ 上结合代数的直和。但 Wedderburn 定理是抽象的，它并不给出同构。

免得你以为这只是无聊的范畴折腾：实数上表示的类似陈述是假的！$\Rep G(\R)$ 是一个合法的范畴，$\hat G$ 也定义在 $\R$ 上，所以看起来可以乐观。但实值不可约表示 $V$ 有三种实性条件（reality conditions）：Schur 引理告诉我们 $\End(V)$ 是 $\R$ 的域扩张上的除代数，于是选项只有 $\R$、$\C$、$\mathbb H$。复表示（$\End(V)\cong\C$）对应 $\hat G$ 的不是 $\R$-点的 $\C$-点；换种说法，特征标不是实值的。即使对实点，也不总有 $\End(V)\cong\R$：若 $\End(V)\cong\mathbb H$，我们显式得到的东西不同构于 $\QCoh(\hat G)$。我们有的结构是：$\Rep G(\R)$ 是可逆的 $\QCoh(\hat G)$-模范畴；但在非代数闭域上，这个模不总是平凡。它某种意义上是 $\QCoh(\hat G)$ 的「范畴线丛」，简称 gerbe。换种说法，$\Rep G(\R)$ 是 $\QCoh(\hat G)$ 的 Brauer 群中的一个元素，就像 $\mathbb H$ 是 $\R$ 的 Brauer 群的非零元素一样。

这与物理中的反常（anomaly）密切相关。$\cC$ 层化到 $M$ 上，但这不意味着它是最简单地层化到 $M$ 上的东西（那个是 $\QCoh(M)$）。它可能是 gerbe，甚至更复杂的东西。随便一个结合代数甚至可能没有非零中心，此时得到的东西离 $\QCoh(Z(A))$ 很远。

\begin{remark}
设 $\Sigma_g$ 是亏格 $g$ 的闭连通可定向曲面。Mednykh 有一个公式，用 $G$ 的表示论计算 $Z_G(\Sigma_g)$，即以自同构加权的 $\Sigma_g$ 上局部系统的计数：
\[
Z(\Sigma_g)=\#\mathrm{Loc}_G(\Sigma_g)=\sum_{\text{irreps }V_\lambda}\left(\frac{|G|}{\dim V_\lambda}\right)^{2g-2}.
\]
这是把数 $Z(\Sigma_g)$ 层化到空间 $\hat G$ 上。一般来说，把 TFT 中的量展布到局部算子代数上，在高维会更有力。
\end{remark}

\srcpage{65}
